\section{Preliminaries}

\subsection{Radial Basis Function (RBF) Neural Networks}
\label{subsec:rbf_nn}

In this paper, Radial Basis Function (RBF) neural networks (NNs) are utilized to approximate the unknown continuous non-linear functions within the robotic manipulator dynamics. 

Generally, an RBF NN can be represented as the linear combination of localized radial basis functions. Given an input vector $x \in \mathbb{R}^m$, the actual output of the RBF NN is expressed in a linearly parameterized form as:
\begin{equation}
	\hat{f}(x) = \hat{W}^T h(x)
\end{equation}
where $\hat{W} \in \mathbb{R}^{m \times p}$ is the weight matrix, and $h(x) = [h_1(x), h_2(x), \dots, h_m(x)]^T \in \mathbb{R}^m$ is the hidden layer activation vector. The Gaussian function is adopted as the basis function for the $i$-th hidden node:
\begin{equation}
	h_i(x) = \exp \left( -\frac{\|x - c_i\|^2}{2b_i^2} \right), \quad i = 1, 2, \dots, n
\end{equation}
where $c_i \in \mathbb{R}^m$ and $b_i > 0$ denote the center vector and the width parameter, respectively.

According to the universal approximation theorem, for any continuous function $f(x)$ defined on a compact set $\Omega_x \subset \mathbb{R}^m$, there exists an ideal constant weight matrix $W^*$ such that:
\begin{equation}
	f(x) = {W^*}^{T} h(x) + \epsilon
\end{equation}
where $\epsilon \in \mathbb{R}^p$ represents the inherent NN approximation error. The ideal weight matrix $W^*$ is defined as:
\begin{equation}
	W^* \triangleq \arg \min_{W} \left[ \sup_{x \in \Omega_x} \| f(x) - W^T h(x) \| \right]
\end{equation}
Furthermore, the approximation error $\epsilon$ is assumed to be bounded on the compact set $\Omega_x$, satisfying:
\begin{equation}
	\|\epsilon\| \le \epsilon_N
\end{equation}
where $\epsilon_N > 0$ is an unknown positive constant.

\begin{remark}
	The linearly parameterized structure of the RBF NN implies that the network output is linear with respect to the weight parameter $\hat{W}$, while maintaining a non-linear mapping via $h(x)$. This characteristic significantly facilitates the derivation of adaptive parameter update laws using the Lyapunov synthesis method, thereby ensuring the closed-loop stability of the robotic system.
\end{remark}

\subsection{Notation}

The dimensions of the principal variables are summarized in Table~\ref{tab:notation}.

\begin{table}[!t]
	\caption{Dimensions of the Main Variables}
	\label{tab:notation}
	\centering
	\begin{tabular}{ll}
		\hline
		Variable & Dimension \\
		\hline
		$e,\dot e$ & $\mathbb{R}^{n\times1}$ \\
		$X=\begin{bmatrix}e\\ \dot e\end{bmatrix}$ & $\mathbb{R}^{2n\times1}$ \\
		$P$ & $\mathbb{R}^{2n\times2n}$ \\
		$B=\begin{bmatrix}0\\M_0^{-1}(q)\end{bmatrix}$ & $\mathbb{R}^{2n\times n}$ \\
		$\varepsilon$ & $\mathbb{R}^{n\times1}$ \\
		$h(X)$ & $\mathbb{R}^{m\times1}$ \\
		$\tilde W$ & $\mathbb{R}^{m\times n}$ \\
		\hline
	\end{tabular}
\end{table}